\documentclass{scrartcl}
\usepackage[T1]{fontenc}
\usepackage[ngerman]{babel}
\usepackage{graphicx}
\usepackage[hidelinks]{hyperref}
\usepackage{subcaption} %erforderlich für subfigure


\begin{document}

	% =========================================
	\section{Zwei Teilabbildungen}
	% =========================================
	
	\begin{figure}[htbp]
		\centering
		
		\begin{subfigure}[t]{0.47\linewidth}
			\centering
			
			\includegraphics[
			width=\linewidth
			]{example-image-a}
			
			\caption{Messung A}
			\label{fig:messung-a}
		\end{subfigure}
		\hfill
		\begin{subfigure}[t]{0.47\linewidth}
			\centering
			
			\includegraphics[
			width=\linewidth
			]{example-image-b}
			
			\caption{Messung B}
			\label{fig:messung-b}
		\end{subfigure}
		
		\caption{Vergleich zweier Messungen}
		
		\label{fig:vergleich}
	\end{figure}
	
	In Abbildung~\ref{fig:vergleich}
	werden zwei Messungen verglichen.
	
	Teilabbildung~\subref{fig:messung-a}
	zeigt die erste Messung.
	
	Teilabbildung~\subref{fig:messung-b}
	zeigt die zweite Messung.

\end{document}
